\section{The Problem: Rights Beyond Reach}

The U.S.\ Constitution promises civilians the right to remain silent, refuse searches, request legal counsel, and walk away when not lawfully detained. But in practice, the system demands that civilians activate these rights with legal precision—under pressure, without training, and at the moment they are least equipped to do so. This failure risks coerced confessions for civilians and ambiguity for officers, rendering rights inaccessible.

Courts have steadily shifted the burden of rights invocation onto civilians. In \textit{Salinas v.\ Texas} (2013) the Supreme Court ruled that silence before arrest could count as evidence of guilt—because the speaker failed to explicitly invoke the Fifth Amendment. In \textit{Berghuis v.\ Thompkins} (2010) a suspect who remained silent for nearly three hours was found to have waived the right to remain silent—simply by answering one question without asserting that right first. In \textit{Doyle v.\ Ohio} (1976) even post‑Miranda silence gains protection only after an explicit invocation. Miranda promises civilians they ``have the right to remain silent''—but omits that explicit invocation is required, misleading civilians. The law offers protection only if the speaker already knows the court‑recognized language—and says it aloud.

Stress turns this legal trap into a human one. People under pressure don’t summon legal formulations; they freeze, err, or speak vaguely—responses courts routinely misread as waiver, guilt, or consent. As cognitive scientist Daniel Willingham explains: ``Thinking is slow, effortful, and uncertain.'' And as Daniel Kahneman warns, ``Even in the absence of time pressure, maintaining a coherent train of thought requires discipline.'' Without a memorized phrase to fall back on, civilians default to silence or improvisation. In a high‑stakes interaction, no one summons legal clarity, and that failure costs them protection.

Existing rights education offers no fix. Current materials rarely provide exact wording, and where example wording is provided, it varies across sources. No single, repeatable script exists—nothing like the clarity Miranda provides to police. Most guides interleave guidance on what to do alongside what to say, and when they address language, they offer multiple suggestions without clarifying what courts actually recognize as invocation. Without a standard phrase to rely on, civilians improvise—and courts treat that as a waiver.

The cost is measurable. A 2004 study of wrongful convictions overturned by DNA evidence found that 29\% involved false confessions—many extracted without legal counsel, amid coercive questioning, or after failed invocations (Drizin \& Leo, \emph{North Carolina Law Review}). Police also must guess whether vague language constitutes invocation. A muttered protest may sound like consent; silence may be misread as defiance. This fuels distrust, undermining safety for all. The lack of shared verbal boundaries makes every encounter harder, tenser, and more dangerous than it needs to be.

The courts demand clarity. But stress, poor education, and legal ambiguity all but guarantee failure—especially for the young, the untrained, and the overpoliced. Constitutional rights remain technically available but functionally unreachable.
